\section{Introduction}
The rapid evolution of automotive architectures has driven a transition from traditional Controller Area Network (CAN) buses to high-speed Automotive Ethernet (AE), enabling bandwidth-intensive applications such as surround-view camera systems, LiDAR data transmission, and over-the-air (OTA) software updates~\cite{matheus2021automotive}. While this architectural shift brings substantial functional advantages, it simultaneously expands the attack surface of in-vehicle networks by exposing vehicles to a broader class of cyber threats previously associated with enterprise and industrial networks, including port scanning, address spoofing, and payload injection~\cite{douss2023state}.

Intrusion detection systems (IDS) have emerged as a primary defense mechanism for in-vehicle networks, and a substantial body of research has been devoted to CAN-based IDS leveraging statistical anomaly detection, machine learning, and deep learning approaches~\cite{lokman2019intrusion}. However, Automotive Ethernet presents fundamentally different traffic characteristics: packets carry variable-length raw byte payloads, traffic is heterogeneous across network segments (e.g., control traffic versus high-bandwidth RTP camera streams), and the protocol stack closely mirrors standard IP networking. These properties render CAN-specific detection methods inapplicable and necessitate purpose-built approaches for Automotive Ethernet.

A central challenge in developing IDS for automotive systems is the scarcity of labeled attack data. Collecting representative attack traffic in real vehicles is operationally constrained, and manually annotating large-scale packet captures is prohibitively expensive. As a result, supervised learning approaches -- which have achieved strong performance in controlled benchmark settings -- are difficult to deploy in practice, as they require comprehensive coverage of attack types at training time and fail to generalize to unseen attack variants~\cite{hindy2020taxonomy}. This motivates an unsupervised formulation of the intrusion detection problem, in which a model is trained exclusively on normal traffic and anomalies are identified as deviations from the learned normal behavior.

Self-supervised learning (SSL) has recently demonstrated remarkable success in learning rich representations from unlabeled data across vision~\cite{he2022masked}, language~\cite{devlin2019bert}, and time-series domains~\cite{yue2022ts2vec}. In particular, masked autoencoding -- wherein a model is trained to reconstruct randomly masked portions of its input -- has proven effective at capturing the underlying generative structure of complex sequential data. This property is well-suited to network traffic anomaly detection: a model that faithfully reconstructs normal packet sequences will produce elevated reconstruction errors when presented with anomalous traffic that deviates from learned patterns.

However, applying masked autoencoding to Automotive Ethernet traffic is non-trivial. Network packets contain two complementary sources of information that exhibit distinct statistical properties: the raw byte content of each packet, which encodes payload structure and protocol fields, and the inter-packet byte differences across a traffic sequence, which capture how each byte position evolves over consecutive packets. Modeling these two information sources with a single homogeneous encoder risks underutilizing either the global sequential structure or the per-byte temporal variation patterns. Furthermore, Automotive Ethernet carries heterogeneous traffic types within the same network, including low-bandwidth control packets and high-bandwidth RTP camera streams, each exhibiting distinct normal behavior and susceptibility to different attack classes.

To address these challenges, we propose a self-supervised dual-stream masked autoencoder for unsupervised intrusion detection in Automotive Ethernet networks. Given a sliding window of 64 consecutive packets, two specialized encoders are trained jointly to exploit the two complementary information sources described above. Global inter-packet structure is captured by a Transformer encoder operating on normalized raw byte content; critically, its positional encoding departs from the fixed sinusoidal scheme common to Transformer architectures and is instead conditioned on the actual inter-packet interval measured at capture time, so that the model reasons about \emph{when} a packet arrived rather than merely its rank within the window, without requiring any upper-layer protocol parsing. Complementing this global view, a residual 1D-CNN treats each byte position as an independent channel and convolves along the packet axis to track fine-grained, per-byte variation across consecutive packets -- variation that a window-wide attention mechanism is comparatively less suited to localize. The two streams are combined through a delta-conditioned gated residual mechanism, in which the CNN stream's own magnitude determines how strongly it corrects the Transformer representation; the fused sequence is then pooled via attention into a single latent vector, from which a conditioned decoder reconstructs the original packet content (Figure~\ref{fig:architecture}). The entire model is trained on normal traffic alone under a masked reconstruction objective with asymmetric loss weighting toward masked positions, and at inference an anomaly is flagged whenever the mean reconstruction error of a window exceeds a threshold calibrated on the normal traffic error distribution.

We evaluate the proposed method on CarDS~\cite{hellemans2025cards}, to our knowledge the first publicly available real-vehicle Automotive Ethernet capture encompassing both control-network and RTP camera-stream traffic across multiple driving scenarios, with attack coverage spanning Nmap-based scanning, IP/MAC address spoofing, and RTP stream replacement. Despite using no attack labels during training, the proposed method reaches an overall F1-score of 0.9808 (AUROC 0.9974, AUPRC 0.9972) across 27 attack files and 23 attack configurations -- performance that, as detailed in Section~\ref{sec:baseline}, closely tracks a fully supervised baseline trained with complete attack labels. Detection is near-perfect on scanning and spoofing attacks (F1 up to 0.9993) and remains competitive on the structurally harder RTP stream replacement category (F1 up to 0.9217), where the attack alters only payload content while leaving packet timing and header structure intact. The architectural ablations reported in Section~\ref{sec:ablation} isolate the source of this performance: both the dual-stream design and the time-aware positional encoding contribute independently, with the latter also markedly extending the range of detection thresholds over which performance stays stable -- an important property for practical deployment (Section~\ref{sec:discussion}).

The main contributions of this paper are as follows:
\begin{itemize}
\item To our knowledge, we present the first anomaly detection study conducted on the CarDS Automotive Ethernet dataset, establishing evaluation protocols and performance baselines for this newly available real-vehicle benchmark.

\item We propose a protocol-agnostic (i.e., parsing-free) dual-stream masked autoencoder for unsupervised Automotive Ethernet intrusion detection: a Transformer encoder over raw packet content, conditioned on a time-aware positional encoding derived from real inter-packet capture timing rather than fixed sequence rank, combined with a residual CNN encoder over inter-packet byte differences through a delta-conditioned gated residual fusion mechanism. The model is trained via packet-level masked reconstruction on unlabeled normal traffic alone, requiring no protocol-specific parsing or labeled attack data. Ablation studies confirm that both the dual-stream design and the time-aware encoding independently improve detection, with the latter also substantially widening the range of thresholds over which performance remains stable.

\item We conduct extensive experiments against representative existing approaches and ablation studies across 27 real-vehicle attack files spanning 23 attack configurations, providing detailed analysis of detection performance across heterogeneous attack categories.
\end{itemize}

The remainder of this paper is organized as follows. Section~\ref{sec:relatedwork} reviews related work on in-vehicle network intrusion detection and self-supervised learning for anomaly detection. Section~\ref{sec:methodology} describes the proposed method in detail. Section~\ref{sec:experiments} presents the experimental setup, results, and ablation studies. Section~\ref{sec:discussion} discusses the broader implications of these results, including the comparison with baseline methods, the role of RTP-specific behavior, and practical deployment considerations. Section~\ref{sec:conclusion} concludes the paper and outlines directions for future work.